%% ============================================================
%% CHAPTER 3 — PHYSICS-INFORMED VULNERABILITY ENGINE (ZONE C)
%% ============================================================

\chapter{Physics-Informed Vulnerability Engine}
\label{ch:physics}

\section{Chapter Overview}
\label{sec:physics_overview}

This chapter develops the three interconnected layers of Zone~C:
(i) the \ac{MC} failure probability engine grounded in \ac{BS7910}:2019
and \acs{IIW} fatigue standards (\cref{sec:physics_mc}),
(ii) the Gulf Coast pipeline network model with physics-derived edge
attributes (\cref{sec:physics_network}), and
(iii) the Bayesian Stackelberg Security Game with LP-enumeration solver
(\cref{sec:physics_game}).  All numerical results in this chapter are
derived from the validated codebase and cross-checked against the
310-test suite described in the Technical Inventory.

\section{PHMSA Empirical Calibration}
\label{sec:physics_phmsa}

\subsection{Dataset}
\label{subsec:physics_phmsa_data}

The probabilistic inputs to the failure model are calibrated against the
\ac{PHMSA} Gas Transmission and Gathering Pipeline Annual Report
dataset~\cite{PHMSA2024}, comprising \num{1996} incident records spanning
2010--2024, with 624 data fields per record.  Filtering for material
and weld-related failure causes yields 216 events with usable seam-type
and material-grade information.

% Figure placeholder
\begin{figure}[htbp]
  \centering
  \includegraphics[width=0.92\textwidth]{fig6_phmsa_cause_distribution}
  \caption{Distribution of incident causes in the PHMSA Gas Transmission
    dataset (2010--2024, $n = \num{1996}$).  Material and weld-related
    failures account for \SI{10.8}{\percent} of all recorded incidents
    and represent the highest-consequence category by median consequence
    cost.}
  \label{fig:phmsa_cause}
\end{figure}

\subsection{Calibrated Distributions}
\label{subsec:physics_calibration}

\Cref{tab:calibrated_scf} presents the \acs{SCF} values and associated
coefficient of variation ($c_v$) calibrated from the PHMSA data for
each seam type.  The lognormal parameterisation follows the standard
practice for material scatter in fracture assessment~\cite{BS7910:2019,
Anderson2005}.

\begin{table}[htbp]
  \centering
  \caption{Calibrated Stress Concentration Factors by Seam Type
    (PHMSA-informed, BS~7910:2019 reference values in parentheses).
    All values are mean $\pm$ one standard deviation.}
  \label{tab:calibrated_scf}
  \begin{tabular}{@{}lcccc@{}}
    \toprule
    Seam Type & $\mu_\mathrm{SCF}$ & $\sigma_\mathrm{SCF}$ &
    $c_v$ & Mean $\Pf$ \\
    \midrule
    Seamless              & 1.00 & 0.05 & 0.050 & 0.53 \\
    DSAW                  & 1.25 & 0.10 & 0.080 & 0.62 \\
    HF-ERW                & 1.50 & 0.15 & 0.100 & 0.66 \\
    LF-ERW (pre-1970)     & 1.80 & 0.20 & 0.111 & 0.74 \\
    Spiral SAW            & 1.40 & 0.12 & 0.086 & 0.64 \\
    Fillet / Lap-welded   & 2.20 & 0.25 & 0.114 & 0.88 \\
    \bottomrule
  \end{tabular}
\end{table}

\section{BS 7910 Failure Assessment Diagram Engine}
\label{sec:physics_fad}

\subsection{Option 1 FAD Curve}
\label{subsec:physics_fad_curve}

The BS~7910:2019 Level~2 Option~1 \ac{FAD} curve is defined by
\cite{BS7910:2019}:
\begin{equation}
  f(\Lr) = \left[1 + 0.5\,\Lr^2\right]^{-1/2}
    \left(0.3 + 0.7\,\mathrm{e}^{-\mu\,\Lr^6}\right),
    \qquad 0 \leq \Lr \leq \Lr^\mathrm{max},
  \label{eq:fad_option1}
\end{equation}
where $\mu = \min(0.001\,E/\sigma_y,\; 0.6)$ and the cut-off
$\Lr^\mathrm{max} = (\sigma_y + \sigma_u)/(2\sigma_y)$.  A flaw is
acceptable if and only if the assessment point $(\Lr, \Kr)$ lies within
the envelope defined by \cref{eq:fad_option1}.

% Figure placeholder
\begin{figure}[htbp]
  \centering
  \includegraphics[width=0.80\textwidth]{fig1_fad_assessment}
  \caption{BS~7910:2019 Level~2 Option~1 Failure Assessment Diagram for
    an API~5L X60 seamless segment.  The nominal assessment point (circle,
    green) lies inside the safe region.  The critical assessment point
    (diamond, red) is outside the FAD envelope, indicating fracture
    instability.  Operating conditions: MAOP at \SI{72}{\percent} SMYS.}
  \label{fig:fad_assessment}
\end{figure}

\subsection{Toughness and Load Ratios}
\label{subsec:physics_kr_lr}

The fracture ratio is $\Kr = K_\mathrm{I} / K_\mathrm{mat}$, where
$K_\mathrm{I}$ is computed from the stress intensity solution for a
surface crack in a cylindrical shell~\cite{BS7910:2019}, and
$K_\mathrm{mat}$ is the characteristic fracture toughness.  The load
ratio $\Lr = \sigma_\mathrm{ref}/\sigma_y$ uses the reference stress
solution for an axial surface flaw~\cite{BS7910:2019,API579:2021}.
Applied stresses assume \ac{MAOP} loading at
\SI{72}{\percent}~\acs{SMYS}, consistent with \acs{ASME}~B31.8:2025
design basis~\cite{ASME_B31_8:2025}.

\section{Monte Carlo Failure Probability Engine}
\label{sec:physics_mc}

\subsection{Simulation Design}
\label{subsec:physics_mc_design}

For each pipeline segment $e$, $N = \num{10000}$ independent realisations
of the material and flaw-geometry random variables are drawn, and the
FAD acceptability criterion is evaluated for each trial:
\begin{equation}
  \Pf(e) = \frac{1}{N}\sum_{k=1}^{N}
    \mathbf{1}\!\left[\Kr^{(k)} > f\!\left(\Lr^{(k)}\right)
    \;\mathrm{or}\;
    \Lr^{(k)} \geq \Lr^\mathrm{max}\right].
  \label{eq:mc_pf}
\end{equation}
Material properties ($\sigma_y$, $E$, $K_\mathrm{mat}$) follow lognormal
distributions with parameters from \cref{tab:calibrated_scf}.  Flaw
dimensions follow truncated normal distributions bounded by API~1104:2021
acceptance criteria~\cite{API1104:2021}.

% Figure placeholder
\begin{figure}[htbp]
  \centering
  \includegraphics[width=0.92\textwidth]{fig4_mc_pf_simulation}
  \caption{Monte Carlo failure probability distributions for four seam
    types ($N = \num{10000}$ trials per segment).  Box-whisker plots
    show the interquartile range; the cross marks the mean.  The ordering
    seamless $<$ DSAW $<$ HF-ERW $<$ LF-ERW matches the empirical failure
    hierarchy in the PHMSA dataset (see \cref{fig:phmsa_cause}).}
  \label{fig:mc_pf}
\end{figure}

\begin{figure}[htbp]
  \centering
  \includegraphics[width=0.88\textwidth]{fig5_pf_scf_sensitivity}
  \caption{Sensitivity of segment failure probability $\Pf$ to stress
    concentration factor (SCF).  A unit increase in SCF from the
    seamless baseline (SCF $= 1.0$) to the fillet-weld value (SCF $= 2.2$)
    amplifies $\Pf$ by a factor of $\approx \num{1.7}$, confirming that
    SCF uncertainty dominates the failure probability variance in this model.}
  \label{fig:pf_scf_sensitivity}
\end{figure}

\section{Gulf Coast Network Model}
\label{sec:physics_network}

\subsection{Topology}
\label{subsec:physics_network_topology}

The synthetic Gulf Coast transmission corridor consists of $|\mathcal{V}|
= 20$ nodes (1 source, 2 compressor stations, 8 delivery points, 3 storage
facilities, 6 junction nodes) and $|\mathcal{E}| = 22$ directed
pipeline segments.  Node positions are sampled from a bounding box
calibrated to the PHMSA Gulf Coast reporting region.  Edges are assigned
diameter, wall thickness, seam type, installation year, and \acs{MAOP} by
stratified sampling from the \ac{PHMSA} fleet distribution.

% Figure placeholder
\begin{figure}[htbp]
  \centering
  \includegraphics[width=0.92\textwidth]{fig11_network_vulnerability_map}
  \caption{Gulf Coast transmission network vulnerability map.
    Edge colour and width encode $\Pf$ (red = high risk).
    Betweenness centrality ranges from $\num{0.00}$ (isolated laterals)
    to $\num{0.26}$ (primary trunk segments).  Segment~SEG\_N012\_N015
    ($\Pf = \num{0.93}$, max betweenness) is the dominant attacker
    target under all three adversary types.}
  \label{fig:network_map}
\end{figure}

\subsection{$\Pf$ Attachment}
\label{subsec:physics_pf_attach}

After topology generation, the \ac{MC} engine is applied to each
segment using its assigned physical attributes, yielding
$\Pf \in [\num{0.29},\; \num{0.93}]$ across the 22 segments.
\Cref{fig:network_map} visualises the resulting vulnerability landscape.

\section{Bayesian Stackelberg Security Game}
\label{sec:physics_game}

\subsection{Game Formulation}
\label{subsec:physics_game_formulation}

The pipeline defender chooses coverage probabilities
$\bmc = (c_1, \ldots, c_n) \in [0,1]^n$ subject to budget
$\sum_i c_i \leq B$.  Each attacker type $\theta \in \Theta$
has prior $p_\theta$ and utility function $\Ua^\theta(i)$ as
defined in \cref{eq:attacker_utility}.  The Bayesian expected
defender utility under coverage $\bmc$, given that the attacker
best-responds to the marginal distribution of defender actions, is:
\begin{equation}
  \Ud(\bmc) = \sum_{\theta \in \Theta} p_\theta \cdot \Ud^\theta(\bmc,
    \mathrm{BR}^\theta(\bmc)),
  \label{eq:bayesian_defender_utility}
\end{equation}
where $\mathrm{BR}^\theta(\bmc) = \argmax_i \bigl[(1 - c_i)\Ua^\theta(i)
\bigr]$ is the attacker's best-response target.

\subsection{LP Formulation and Strong Stackelberg Equilibrium}
\label{subsec:physics_lp}

Following Paruchuri et al.~\cite{Paruchuri2008}, for each candidate
target $q$, the \ac{SSE} is found by solving:
\begin{align}
  \min_{\bmc} \quad & -(1-\delta)\,\Ua(q)\,c_q
    \label{eq:lp_obj} \\
  \text{s.t.} \quad
    & \sum_{i=1}^n c_i \leq B, \label{eq:lp_budget} \\
    & c_i \in [0,1] \;\forall i, \label{eq:lp_bounds} \\
    & \eps\,\Ua(q)\,c_q - \eps\,\Ua(j)\,c_j \leq \Ua(q) - \Ua(j)
      \;\forall j \neq q, \label{eq:lp_br}
\end{align}
where $\delta = \num{0.25}$ is the protection factor (reduction in
attacker utility when a segment is covered) and $\eps = \num{1e-4}$
enforces the strict-best-response constraint~\eqref{eq:lp_br}.  The
correct sign convention for constraint~\eqref{eq:lp_br} --- $+\eps$
on the $q$-term and $-\eps$ on the $j$-term --- is essential and was
verified against the full test suite.

\subsection{Attacker Type Profiles}
\label{subsec:physics_attacker_types}

\begin{table}[htbp]
  \centering
  \caption{Attacker type profiles used in the Bayesian SSE.
    Weight parameters $w_{\Pf}$, $w_v$, $w_b$ define how each type
    computes its target utility (Equation~\ref{eq:attacker_utility}).}
  \label{tab:attacker_profiles}
  \begin{tabular}{@{}lcccccc@{}}
    \toprule
    Type & Prior $p_\theta$ & $w_{\Pf}$ & $w_v$ & $w_b$
      & Description \\
    \midrule
    \textsc{Strategic}    & 0.50 & 0.6 & 1.0 & 0.0
      & Maximises asset value damage \\
    \textsc{Opportunistic}& 0.30 & 1.0 & 0.2 & 0.0
      & Targets structurally weakest pipe \\
    \textsc{State Actor}  & 0.20 & 0.5 & 0.7 & 1.0
      & Maximises network disruption \\
    \bottomrule
  \end{tabular}
\end{table}

\subsection{Game Results}
\label{subsec:physics_game_results}

At a budget fraction of \SI{30}{\percent} ($B = \num{6.6}$ coverage
units over 22 segments), the Bayesian \ac{SSE} yields:

\begin{itemize}
  \item Defender utility: $\Ud^* = \num{-0.477}$
  \item Coverage effectiveness versus zero-budget baseline: $+\SI{52.3}{\percent}$
  \item Risk reduction versus uniform allocation: $\SI{17.0}{\percent}$
  \item Budget utilisation: $\num{6.6}/\num{6.6}$ (\SI{100}{\percent})
  \item Primary defended segment: SEG\_N012\_N015
    ($\Pf = \num{0.93}$, maximum betweenness centrality)
\end{itemize}

% Figure placeholders
\begin{figure}[htbp]
  \centering
  \includegraphics[width=0.92\textwidth]{fig13_stackelberg_coverage_map}
  \caption{Bayesian Stackelberg SSE coverage probabilities $c_i$ over the
    Gulf Coast network.  Edge width encodes $c_i$; node size encodes
    betweenness centrality.  The physics-informed allocation concentrates
    resources on trunk segments with both high $\Pf$ and high betweenness,
    in contrast to the uniform baseline ($c_i = \num{0.30}$ for all $i$).}
  \label{fig:stackelberg_coverage}
\end{figure}

\begin{figure}[htbp]
  \centering
  \includegraphics[width=0.92\textwidth]{fig14_budget_sensitivity}
  \caption{Budget sensitivity analysis across 18 budget levels
    ($B/N \in [\num{0.05}, \num{0.90}]$).
    \textit{Left:} Defender and attacker utilities.
    \textit{Centre:} Coverage effectiveness (\% improvement vs.\ baseline).
    \textit{Right:} Marginal return $\partial \Ud / \partial B$.
    Diminishing returns are evident beyond $B/N \approx \num{0.60}$,
    providing an evidence base for the \SI{30}{\percent} budget choice.}
  \label{fig:budget_sensitivity}
\end{figure}
